\scanpage{77}
the $n$.th power homomorphism in the Picard prescheme is of finite type.

I tell you about this, namely (i), because of the method of proof, involving of course considerations of non flat descent. The fact that I do not have any good effectivity criterion does not hamper, by just recalling what the effectivity of a given descent datum means. \deleted{By a} Now it turns out that by a slightly more careful analysis of the situation, one can prove the following theorem, of a type very close to the one you have proved recently, and \inserted{\handwritten{to}} some you still want to prove as I understand it\deleted{\unsure{:}}.

\uline{Theorem} Let $S$ be an integral noetherian scheme, $X$ and $X\handwritten{'}$ proper over $S$ and $f:X\handwritten{'\to}X$ a surjective $S$-morphism, look at the corresponding homomorphisme for the Picard functors $f^{\bullet}:\underline{\underline{\unsure{\mathbb{P}|P}\mathrm{ic}}}_{X/S}\handwritten{\to}\text{\uuline{Pic}}_{X'/S}$ . Assume: a) the existence problem\deleted{s} A \deleted{and B} defined below for $X/S$ ha\deleted{v}\handwritten{s}\deleted{e} always a solution (this is certainly true when $X/S$ is projective). b) the morphism\deleted{\unsure{e}} $f_{s}:X\handwritten{'}_{s}\handwritten{\to}X_{s}$ induced on the generic fiber is a morphism of descent, i.e. $\underline{O}_{X_{s}}\handwritten{\to}\deleted{\unsure{g}}\handwritten{f_{*}}(\underline{O}_{X\handwritten{'}_{s}})\handwritten{\unsure{\rightrightarrows|\substack{\longrightarrow\\[-0.7ex]\rule{1.4em}{0.4pt}}}}h_{\handwritten{\unsure{*|\uparrow}}}(\underline{O}_{X''_{s}})$ is exact. Then, \uline{provided we replace $S$ by a suitable non empty open set},\margin{\handwritten{\unsure{$\epsilon$|e|c|$\in$} not\\ needed ?}} the homomorphism\deleted{\unsure{e}} $f^{\bullet}$ is representable by \deleted{a} quasi-affine morphism\deleted{s}, more sp\unsure{e}cifically \deleted{the} in the factorisation of $f^{\circ}$via the functor representing suitable descent data, $f^{\bullet}=vu$ with $u$\unsure{.} affine \deleted{\unsure{as you know, both}} and $v$ a monomorphism (as you well know), $v$ is in fact representable by \inserted{\handwritten{finite direct sums of}} immersions.

\uline{Corollary \handwritten{1}} Without assuming b), but instead in a) allowing $X/S$ to be replaced by suitable other schemes $X_{\unsure{i|1}}$ finite over $\unsure{\mathbb{X}|X}$, the same conclusion holds, namely $f^{\circ}$ is representable by \deleted{\unsure{finite|finit}} quasi-affine morphisms.

This follows from the theorem, using a suitable factorisation of
% archive sheet number: 76
